\question[3]
In einer Speicherzelle von Johnny steht der Wert \texttt{02.015}.
Erläutere, ob es sich dabei um eine Zahl oder um einen Befehl handelt.
Gehe dabei auf das Von-Neumann-Prinzip ein, nach dem Programm und Daten im selben Speicher liegen.

\begin{solutionbox}{5cm}
Das lässt sich am Wert allein nicht entscheiden. Beim Von-Neumann-Rechner liegen
Programm und Daten im selben Speicher, und einer Speicherzelle sieht man nicht an,
was sie enthält.

Wird die Zelle über den Befehlszähler geholt und ins Befehlsregister übernommen,
wird sie als Befehl \texttt{ADD 15} ausgeführt. Wird sie dagegen von einem Befehl wie
\texttt{TAKE} gelesen, wird sie als Zahl (2015) verwendet.
\end{solutionbox}
